\section{Conclusion}
\label{sec:conclusion}

We compared a particle swarm, random search and the default settings for tuning a small waste detector, giving each search the same budget of 32 proxy trainings.

On RQ1, PSO did not beat random search in any way that matters. In \genone (studio data) the three arms were statistically tied at 0.676, 0.672 and 0.668 \mapfifty. In \gentwo both tuned arms beat the defaults on the mixed test set (random search by $+0.017$, PSO by $+0.020$), but PSO and random search could not be told apart on the mixed, studio or real-world test sets. Their only significant difference, $+0.003$ on the Indian set, is on scores near 0.02 and has no practical meaning. On RQ2, the short proxy did rank settings like longer training ($\rho = 0.976$, $n = 8$, \genone only), but the small gaps it showed between the best settings disappeared after full training. On RQ3, the answer changes with the domain. The tuning gains did not reach real-world photographs (0.298 to 0.310), and all arms failed on the Indian DWSD photographs (0.023 to 0.027), partly because of the merged boxes in DWSD. Real-world training data had a much larger effect: on the same images it raised real-world \mapfifty from 0.038--0.096 (\genone) to 0.338--0.353 (\gentwo), at a studio cost of 0.074 to 0.105.

We take two things from this. A claim that a metaheuristic tunes a model better should include the defaults and random search at an equal budget, and confidence intervals on a test split that is used once. Without the random-search arm we would have credited our \gentwo gain to the swarm. And for waste detection in real scenes, especially Indian ones, the training data matter far more than the optimizer. Our next steps are Indian training data with a box for each item and an ensemble of the three models (Section~\ref{sec:future}).
